\documentclass[letterpaper,10pt]{article}

\usepackage{latexsym}
\usepackage[empty]{fullpage}
\usepackage{titlesec}
\usepackage{marvosym}
\usepackage[usenames,dvipsnames]{color}
\usepackage{verbatim}
\usepackage{enumitem}
\usepackage[hidelinks]{hyperref}
\usepackage{fancyhdr}
\usepackage[english]{babel}
\usepackage{tabularx}
\usepackage{multicol}
\input{glyphtounicode}

\usepackage[default]{sourcesanspro}
\usepackage[T1]{fontenc}

\pagestyle{fancy}
\fancyhf{} 
\fancyfoot{}
\renewcommand{\headrulewidth}{0pt}
\renewcommand{\footrulewidth}{0pt}

\setlength{\footskip}{3.7 pt}
\addtolength{\oddsidemargin}{-0.5in}
\addtolength{\evensidemargin}{-0.5in}
\addtolength{\textwidth}{1in}
\addtolength{\topmargin}{-.5in}
\addtolength{\textheight}{1.0in}

\urlstyle{same}

\raggedbottom
\raggedright
\setlength{\tabcolsep}{0in}

\titleformat{\section}{
  \vspace{-4pt}\centering
}{}{0em}{}[\color{black}\titlerule\vspace{-5pt}]


\pdfgentounicode=1

\newcommand{\resumeItem}[1]{
  \item\small{
    {#1 \vspace{-2pt}}
  }
}

\newcommand{\resumeSubheading}[4]{
  \vspace{-2pt}\item
    \begin{tabular*}{0.97\textwidth}[t]{l@{\extracolsep{\fill}}r}
      \textbf{#1} & #2 \\
      \textit{\small#3} & \textit{\small #4} \\
    \end{tabular*}\vspace{-7pt}
}

\newcommand{\resumeSubSubheading}[2]{
    \item
    \begin{tabular*}{0.97\textwidth}{l@{\extracolsep{\fill}}r}
      \textit{\small#1} & \textit{\small #2} \\
    \end{tabular*}\vspace{-7pt}
}

\newcommand{\resumeProjectHeading}[2]{
    \item
    \begin{tabular*}{0.97\textwidth}{l@{\extracolsep{\fill}}r}
      \small#1 & #2 \\
    \end{tabular*}\vspace{-7pt}
}

\newcommand{\resumeSubItem}[1]{\resumeItem{#1}\vspace{-4pt}}

\renewcommand\labelitemii{$\vcenter{\hbox{\tiny$\bullet$}}$}

\newcommand{\resumeSubHeadingListStart}{\begin{itemize}[leftmargin=0.15in, label={}]}
\newcommand{\resumeSubHeadingListEnd}{\end{itemize}}
\newcommand{\resumeItemListStart}{\begin{itemize}}
\newcommand{\resumeItemListEnd}{\end{itemize}\vspace{-5pt}}

\begin{document}



\begin{center}
    {\LARGE Zhijie Lyu} \\ \vspace{2pt}
    \begin{multicols}{2}
    \begin{flushleft}
    %\href{{your github page link}}{my github}\\
    {Email \\ }
    \href{zhijie.lyu@gmail.com}
    {zhijie.lyu@gmail.com} 
    \end{flushleft}
    \begin{flushright}
    \href{https://www.linkedin.com/in/zhijie-lyu-7b461b293/}{https://www.linkedin.com/in/zhijie-lyu-7b461b293/}
    \href{https://github.com/alsymiya}{https://github.com/alsymiya}
    %{Phone: (1)631-202-9630}
    \end{flushright}
    \end{multicols}
\end{center}


%-----------EDUCATION-----------
\vspace{-2pt}
\section{Education}
  \resumeSubHeadingListStart
    \resumeSubheading
      {Stony Brook University -- SUNY}{January 2020 -- May 2026}
      {PhD student in Mechanical Engineering with a minor in Computer Science}{Stony Brook, New York, USA}
  
      \resumeSubheading
      {Stony Brook University -- SUNY}{August 2018 -- January 2020}
      {Master's student in Mechanical Engineering, Manufacturing Direction}{Stony Brook, New York, USA}

    \resumeSubheading
      {Southwest Jiaotong University}{August 2014 -- May 2018}
      {Bachelor in Mechanical Engineering, Mechatronic Direction}{Sichuan, Chengdu, PRC}

  \resumeSubHeadingListEnd


%-----------EXPERIENCE-----------
\section{Experience}
  \resumeSubHeadingListStart
    \resumeSubheading
      {Coursework}{August 2018 -- December 2021}
      {GPA during master's: 3.93; GPA during Ph.D.: 3.80}{}
      \resumeItemListStart
        \resumeItem{Topics: Robotics, Computer-Aided Design, Data Driven Design, Analysis of Algorithms, Computational Geometry}
      \resumeItemListEnd
    
    \resumeSubheading
      {Ph.D. Research}{January 2020 -- Present}
      {Focus: Data-Driven Design, Data Structures, and Algebraic Graph Theory}{}
      \resumeItemListStart
        \resumeItem{\textbf{Journal Publication Record}: $1$ for Data-Driven Synthesis Pipeline (Selected as featured article for ASME JMD); $1$ for Mechanism Dataset; $3$ for Efficient Geometric Constraint Solving; $1$ for Mechanism Cognate Generation, as second author. } 
        \resumeItem{\textbf{Journal Review Record}: multiple manuscripts from ASME-JMD, ASME-JMR and IEEE-THMS \\}
        \resumeItem{\textbf{On-going Research:} Differentiable and parallelized computation framework for CPUs and GPUs. }
      \resumeItemListEnd

    \resumeSubheading
    {Teaching}{January 2020 -- Present}{Served as a teaching assistant for undergraduate students and high school students}{}
    \resumeItemListStart
        \resumeItem{Worked as TA for four undergraduate courses: Intro to M. Eng., Statics, Dynamics, Numerical Methods, and Mechanism Design}
        \resumeItem{Guided a high school student to develop a framework for mechanism design; got the work published in MDPI}
      \resumeItemListEnd
    
    \resumeSubheading
    {Academic Presentation}{August 2022, 2023, 2024, and 2025}
    {Participated in four ASME IDETC/CIE and presented our academic research results. }{}

    \resumeSubheading
    {Software Development (Funded by NSF and Mechanismic Inc.)}{May 2019 -- Present}
    {Development and Evaluation of \href{motiongen.io}{MotionGen} (\textbf{motiongen.io}, a web-based CAD application)}{}
    \resumeItemListStart
        \resumeItem{\textbf{Display}: implementation of customizable shape elements on the canvas using paper.js \textbf{(2019 -- 2020)}}
        \resumeItem{\textbf{Mechanism Simulation}: \begin{enumerate}
            \item Cooperated on a unified, generalizable, differentiable and efficient planar linkage simulator using approximation \textbf{(2022)}
            \item Extended the simulation capability using a Pebble Game framework for gears and wheels into the linkage system with precise constraint system. \textbf{(2024 -- 2025)}
        \end{enumerate}} 
        \resumeItem{\textbf{Mechanism Synthesis}:  \begin{enumerate}
            \item Deployed the G-manifold algebraic fitting method for motion synthesis \textbf{(2022)}
            \item Deployed a backend server using Google Cloud and a neural network pipeline \textbf{(2023)}
            \item Designed a compact data structure that reduced the storage of the same information data to $1/100$ of its original in the code base, and reduced the query time with the use of customized topology dictionary \textbf{(2024)}
            \item Generated a massive mechanism dataset for these mechanisms \textbf{(2024)}
            \item Added differential evolution method to the synthesis framework \textbf{(2025)} 
        \end{enumerate}} 
        \resumeItem{\textbf{On-going Development (2025-2026)}: 
        \begin{enumerate}
            \item Vectorization (matrix-formatted parallel programming) of the simulator for MotionGen, intended for real-time synthesis
            \item Differentiable physics intended for synthesis fine-tuning 
            \item Multibody dynamics, intended for real-time rendering 
        \end{enumerate} }
        \resumeItem{\textbf{Impact}: The number of registered users is approximately $100,000$ as of November 2025, since it had a registration system in 2020. A sharp increase in user registration occurred in August 2023 when a popular mechanism design channel Engineezy used our application to design a mechanism prototype. } 
        \resumeItemListEnd

    

    
  \resumeSubHeadingListEnd

  %-----------PROGRAMMING SKILLS-----------
\section{Technical Skills, Language Skills, and Interests}
 \begin{itemize}[leftmargin=0.15in, label={}]
    \small{\item{
     \textbf{Programming Languages}{: JavaScript, Python, MatLab scripts, C++ (OOP)} \\
     
     \textbf{Writing}{: \LaTeX, Office} \\
     \textbf{Languages}{: Cantonese (native), Mandarin Chinese (native), English (fluent)} \\
     \textbf{Interests}{: Prompt Engineering, Software Development, Applied Math Science, High Performance Computation (HPC), Document Translation, Speedrunning in video games} \\ 
    }}
 \end{itemize}


%-----------PUBLICATIONS-----------
\newpage
\section{Selected Publications (The first three have been implemented in MotionGen 2024)}
 \begin{itemize}[leftmargin=0.15in, label={}]
    \item{
    \textbf{A Unified Real-Time Motion Generation Algorithm for Approximate Position Analysis of Planar N-Bar Mechanisms - ASME Journal of Mechanical Design, June 2024, 146(6): 063302 (12 pages)}{\\ Authors: \textbf{ \emph{Zhijie Lyu}}, Wei Liao, Anurag Purwar}
    {\\ \small{\begin{itemize}
        \item Extended the prismatic joint approximation trick from specific cases in textbook to arbitrary combination cases; 
        \item Defined a standardized matrix-set data structure for kinematically simulated samples, and their structural complexity analysis; 
        \item \emph{This documents the core of the 2022-2025 MotionGen simulator.}
    \end{itemize} 
    }}
    }
    \item{\textbf{Deep Learning Conceptual Design of Sit-to-Stand Parallel Motion Six-Bar Mechanisms - ASME Journal of Mechanical Design, July 2024, 147(1): 1-18}{\\ Authors: \textbf{\emph{Zhijie Lyu}}, Anurag Purwar}
    {\\ \small{\begin{itemize}
        \item Improved an atlas-based and data-driven approach for planar linkage path synthesis during data preprocessing; 
        \item Used a Variational Auto-Encoder (VAE) along with three metrics to rank solution candidates retrieved from the dataset. 
        \item \emph{The work is selected as a featured article on ASME Journal of Mechanical Design by editor Faez Ahmed from MIT. }
    \end{itemize}}}
    \item{\textbf{A Dataset of 3M Single-DOF Planar 4-, 6-, and 8-bar Linkage Mechanisms with Open and Closed Coupler Curves for Machine Learning-Driven Path Synthesis, ASME Journal of Mechanical Design (Special Issue for IDETC-CIE 2024)}{\\ Authors: Anar Nurizada, Rohit Dhaipule, \emph{\textbf{Zhijie Lyu}}, Anurag Purwar}
        {\\ \small{\begin{itemize}
        \item A dataset of three million samples of R-joint and/or P-joint mechanisms is created with the simulation package of MotionGen; 
        \item \emph{This dataset, together with a compact kinematic structure library, is implemented onto MotionGen using Google Cloud API with the synthesis method presented in the second publication. }
    \end{itemize}}}
    }
    \item{\textbf{Point-based Models: A Unified Approach for Geometric Constraint Analysis of Planar N-Bar Mechanisms with Rotary, Sliding, and Rolling Joints, ASME Journal of Mechanisms and Robotics (Special Issue for IDETC-CIE 2024)}{\\ Authors: \textbf{\emph{Zhijie Lyu}}, Anurag Purwar}
        {\\ \small{
        \begin{itemize}
        \item Developed the methodology and data structure to extend the constraint graph approach from two types of constraints (distance, angle) to three types (angle-angle); This is intended for the implementation of wheels/gears into the linkage system. 
        \item \emph{This is theory part that says wheels can be represented as linkages, so we can analyze all the elements like analyzing a truss structure. }
        \end{itemize}
        }}
    }
    \item{\textbf{Construction, Reduction, and Decomposition of Planar Mechanisms via an Extended Pebble Game Framework, ASME Journal of Mechanisms and Robotics, December 2025, 17(12): 121009 (15 pages) } {\\Authors: \textbf{\emph{Zhijie Lyu}}, Anurag Purwar}
        {\\ \small{\begin{itemize}
        \item Developed a modified pebble game framework to extend the constraint graph approach from two types of constraints (distance, angle) to three types (angle-angle); The framework analyzes the planar geared linkage system in polynomial time. 
        \item The package has been completed in a Python-based prototype. 
        \item \emph{This is a polynomial-time cost algorithm for general mobility analysis, but is expanded from linkages to gears and linkages. }
    \end{itemize}}}
    \item{\textbf{Path Generative Model Based on Conditional $\beta$-Variational Auto Encoder for Four-Bar Mechanism Design, ASME Journal of Mechanisms and Robotics,  June 2025, 17(6): 061004}{\\ Authors: Anar Nurizada, \textbf{ \emph{Zhijie Lyu}}, Anurag Purwar}
    {\\ \small{\begin{itemize}
        \item The paper itself is mostly led by Anar Nurizada. The research goal is given some paths, the neural network outputs an initial guess on the mechanism configurations. 
        \item I made an automated script on generating the four-bar mechanism cognates for any given four-bar linkage, and contributed to the text related to this part. 
    \end{itemize}}}
    
    \item{\textbf{(Submitted) Data-Driven Design using Differentiable Physics: Sampling Density for Reliable Initialization into Basin of Attraction, ASME IDETC/CIE, 2026}{\\ Authors: \textbf{ \emph{Zhijie Lyu}}, Suhas Gangreiddy, Anurag Purwar}
    {\\ \small{\begin{itemize}
        \item Integrated machine learning with the customized differentiable simulator to generate precise and satisfactory designs. 
        \item Experimented on different topologies to check their respective sampling densities. 
    \end{itemize}}}
 \end{itemize}



%-----------CERTIFICATIONS-----------
\begin{comment}
\section{Course and/or Hobby-based Projects}
 \begin{itemize}[leftmargin=0.15in, label={}]
   \small{
    \item{\textbf{Stony Brook University - MEC529 - Robotics}: \\
    Baxter Robot Arm Writing a Letter on Vertical Board (MatLab script)
    \item{\textbf{Stony Brook University - MEC572 - Computer Aided Design}: \\ 
    Interactive Manipulation of NURBS Surfaces (C++ and JavaScript) }

    \item{\textbf{Stony Brook University - MEC596 - Projects in Mechanical Engineering}: \\ Bluetooth-based connection between MatLab script and Arduino. Implementation of Q-learning algorithm for robot car clash prevention. }

    \item{\textbf{Stony Brook University - AMS561 - Introduction to Computational and Data Science}: \\ Sudoku Solver using Depth-first search for 3x3 and 4x4. (Python)
    }

    \item{\textbf{Internship - Online Assessment for a Software Development Engineer role}: \\
    Deployed a Weather Forecast App on Heroku with Create, Read, Update, and Delete functions (Python Flask, JavaScript, HTML, and CSS)}
    }

    
        \item{Hobby - Novel / Video Translation: \\ Was an active Chinese-Japanese/English translator for Japanese novels and English video guides for competitive surf}
    
    
}   
\end{itemize}
\end{comment}
%-----------PROGRAMMING SKILLS-----------

\begin{comment}
    \section{Professional Activities}
\begin{itemize}[leftmargin=0.15in, label={}]
\item \textbf{Literature Review Service}{: 
\begin{itemize}
    \item Reviewed two submissions as an invited reviewer of the ASME Journal of Mechanical Design (JMD)
\end{itemize}
} \\ 
\item \textbf{Teaching Assistant}: {
\begin{itemize}
    \item MEC 101 - Freshman Design Innovation (2021 Fall, 2024 Fall) 
    \item MEC 260 - Engineering Statics (2021 Spring)
    \item MEC 262 - Engineering Dynamics (2025 Spring)
    \item MEC 410 - Design of Machine Elements (2022 Spring)
\end{itemize}}
\end{itemize}
\end{comment}



\end{document}
